% BEGIN EXPONENT-COMPRESSION SUPPLEMENT 2026-09-06
\ifdefined\XCStandalone
\else
\clearpage
\subsection*{Research continuation: exponent compression and overlapping obstructions}
\fi
\noindent\textbf{Author: ChatGPT. Status: ordinary mathematical proofs with explicitly
cited external inputs. This supplement is not a proof or disproof of standard
abc. Its new Lean companion has not been compiled in the present execution.}

\section{The two remaining difficulties and the notation}
\label{sec:xc-scope}
The preceding moving-support supplement controls small boundary primes only
under restrictions on the number and size of the prime factors of the
auxiliary norm. We give a different compression mechanism, an elementary
local estimate on exact-power representations, and an obstruction showing
that large boundary multiplicities and prime auxiliary norms occur together.
The results below are new derivations in this project; no claim of priority
within the literature is made.

Let $a,b,c$ be positive integers with $a+b=c$ and $\gcd(a,b)=1$, and put
\[
 T=abc,\quad R=\rad(T),\quad t=\log c,\quad
 z=a+b\zeta,\quad M=N(z)=a^2+ab+b^2,
 \qquad\zeta^2-\zeta+1=0.
\]
Write $\mathcal E=\Z[\zeta]$. For $w=x+y\zeta$ define
\[
 P(w)=xy(x+y),\quad H(w)=\max(|x|,|y|,|x+y|),
 \quad N(w)=x^2+xy+y^2.
\]
The multiplication, conjugation, Euclidean factorization and unit rotations
are those of the preceding Eisenstein supplement. In particular, for
primitive $w$,
\begin{equation}\label{eq:xc-geometric}
 \gcd(N(w),|P(w)|)=1,\qquad
 \frac34H(w)^2\le N(w)\le H(w)^2,\qquad
 |P(w)|\le\frac{2}{3\sqrt3}N(w)^{3/2}.
\end{equation}
The last inequality follows from $4|P|\le H^3$. A primitive zero-boundary
pair is a unit. All radical statements below concern nonzero boundaries.
We retain the exact excess and credit
\[
 E_3(T)=\prod_{p\mid T}p^{(v_p(T)-3)_+},\qquad
 C_3(T)=\prod_{p\mid T}p^{(3-v_p(T))_+},\qquad
 TC_3(T)=R^3E_3(T).
\]
For real $Y\ge2$ let $S_Y(T)=\prod_{p\mid T,\ p\le Y}p^{v_p(T)}$ and
$E_{3,>Y}(T)=\prod_{p\mid T,\ p>Y}p^{(v_p(T)-3)_+}$.

\section{Compressing arbitrarily many norm primes to one power}
\label{sec:xc-content}
Use the actual factorization
\begin{equation}\label{eq:xc-factorization}
 z=u(1+\zeta)^\epsilon\prod_{i=1}^r\pi_i^{e_i},\qquad
 \epsilon\in\{0,1\},\quad N(\pi_i)=q_i\equiv1\pmod3.
\end{equation}
The $q_i$ are distinct and only one conjugate orientation occurs for each.
The case $r=0$ gives $(1,1,2)$ and is henceforth handled separately.

\begin{definition}[Common exponent content]\label{def:xc-content}
Set $g=\gcd(e_1,\ldots,e_r)$ and
$w=\prod_i\pi_i^{e_i/g}$. Put $Q=N(w)$. Thus
\begin{equation}\label{eq:xc-content}
 z=u(1+\zeta)^\epsilon w^g,\qquad
 N(z)=3^\epsilon Q^g,\qquad Q\ge7,\quad\gcd(Q,6)=1.
\end{equation}
Every positive common divisor of the $e_i$ could be used instead of their gcd.
\end{definition}

\begin{lemma}\label{lem:xc-primitive}
The pair $w$ and every power $w^n$, $n\ge1$, are primitive. Neither $P(w^n)$
nor $w^{3n}+\bar w^{3n}$ is zero. Furthermore
\[
 \frac g2\log Q\le t,\qquad \log Q\le\frac{2t}{g}.
\]
\end{lemma}
\begin{proof}
A common rational prime divisor of the two coordinates would be an inert
prime factor, two ramified factors, or both orientations of a split prime.
None occurs in $w^n$. A primitive element with $P=0$ is a unit, inconsistent
with norm $Q^n>1$. The identity
\begin{equation}\label{eq:xc-trace}
 w^3+\bar w^3=(x-y)(2x+y)(x+2y)
\end{equation}
shows that a primitive pair with this trace zero has norm three: the three
possible coordinate ratios give $(1,1)$, $(1,-2)$ or $(-2,1)$ up to signs.
This contradicts $3\nmid Q$. Finally $N(z)\le c^2$ proves both inequalities.
\end{proof}

The only transcendence input in the next theorem is the fixed-degree,
fixed-number-of-logarithms estimate (1.2) of Bugeaud's Theorem~1.1
\cite{moving:bugeaud}. In degree two and for at most two logarithms, it gives
$\log|\Lambda|\ge-C\prod h^*(\alpha_i)\log B$ with an effective absolute
constant. No algebraic generator is required to be fixed across triples.

\begin{theorem}[Support-free compressed angular estimate]\label{thm:xc-angle}
There is an effective absolute constant $C>0$ such that every representation
\eqref{eq:xc-content} satisfies
\begin{equation}\label{eq:xc-angle}
 \left|\left(\frac z{\bar z}\right)^3-1\right|
 \ge \exp\{-C\log Q\log(3g)\}.
\end{equation}
Consequently its angular penalty divided by $\log c$ is at most
$2C\log(3g)/g$, uniformly in the number and sizes of all norm primes.
\end{theorem}
\begin{proof}
Set $h=(w/\bar w)^3$. Then $h^*(h)\le3\log Q$, since each conjugate of
$w$ has absolute value $\sqrt Q$ and the standard height laws give
$h(h)\le6h(w)=3\log Q$. Also
\[
 \lambda_z=(z/\bar z)^3=(-1)^\epsilon h^g.
\]
A unit contributes its sixth power, and
$(1+\zeta)/(1+\bar\zeta)=\zeta$. The cubic difference identity
$z^3-\bar z^3=3\sqrt{-3}P(z)$ proves $\lambda_z\ne1$.
Choose $\log h=i\theta$, $|\theta|\le\pi$, and an integer $j$ such that
\[
 \Lambda=g\log h+(\epsilon-2j)\log(-1)
\]
has imaginary part in $[-\pi,\pi]$. It is nonzero and
$|\epsilon-2j|\le g+1$, so its coefficient bound is at most $3g$.
The cited estimate, followed by $|e^{iv}-1|\ge2|v|/\pi$ for $|v|\le\pi$,
gives \eqref{eq:xc-angle} after absorbing a fixed additive constant.
If a coefficient vanishes, delete it or reorder the remaining terms as
required by the source theorem. Lemma~\ref{lem:xc-primitive} gives the final
uniform ratio. This proves a quantitative estimate, not the assertion that
$g$ tends to infinity on all abc triples.
\end{proof}

\section{An elementary bound for local multiplicities}
\label{sec:xc-local}
The exact-power representation permits a bound linear in $p$, using order
and binomial expansion instead of a general $p$-adic logarithmic estimate.
Here $v_p$ on rational integers has its usual normalization. At a place of
$\Q(\sqrt{-3})$ we use the extension satisfying $v(p)=1$.

\begin{theorem}[Local bounds with all exceptional primes included]
\label{thm:xc-local}
For \eqref{eq:xc-content}, and every prime $p\ge5$ dividing $T$, one has
\begin{equation}\label{eq:xc-local-large}
 v_p(T)\log p\le\frac{p+1}{2}\log Q+v_p(g)\log p.
\end{equation}
At two and three one has
\begin{align}
 v_2(T)\log2&\le3\log Q+v_2(g)\log2,\label{eq:xc-local-two}\\
 v_3(T)\log3&\le\tfrac32\log Q+v_3(g)\log3.
 \label{eq:xc-local-three}
\end{align}
When $\epsilon=1$, in fact $v_3(T)=0$. When $\epsilon=0$,
$v_3(T)=v_3(P(w))+v_3(g)$.
\end{theorem}
\begin{proof}
First suppose $p\ge5$ and $p\mid T$. By \eqref{eq:xc-geometric}, $p\nmid Q$.
In the quadratic algebra modulo $p$, conjugation is trivial on the split
components when $p\equiv1\pmod3$ and is Frobenius when $p\equiv2\pmod3$.
Write $\chi=1$ or $-1$ respectively. The element $w/\bar w$ has norm one,
so its order divides $p-\chi$. Since $3\mid p-\chi$, the order $d$ of
$h=(w/\bar w)^3$ divides $(p-\chi)/3$, and $p\nmid d$.
The numerator of $h^d-1$ is $3\sqrt{-3}P(w^d)$, hence
\[
 s:=v(h^d-1)=v_p(P(w^d))\ge1.
\]
This valuation is the same at both components in the split case, because
$P(w^d)$ is a rational integer and all other factors are units.

For an odd prime, binomial expansion gives
$v((1+p^s u)^n-1)=s+v_p(n)$ when $u$ is a unit and $s\ge1$:
a multiplier prime to $p$ preserves the leading valuation, and raising to
$p$ increases it by exactly one. For the latter assertion the first term
has valuation $s+1$ and all the other terms have valuation at least $s+2$.
This proof also works in each split component.
If $\epsilon=0$, $d\mid g$ and $v_p(T)=s+v_p(g)$. If $\epsilon=1$,
$h^g=-1$ modulo $p$, so $d\mid2g$ and $h^g-1$ is a unit. Factoring
$h^{2g}-1$ gives the same formula $v_p(T)=s+v_p(g)$.
Finally
\[
 s\log p\le\log|P(w^d)|<\tfrac32d\log Q
 \le\tfrac12(p+1)\log Q,
\]
using \eqref{eq:xc-geometric} and Lemma~\ref{lem:xc-primitive}.

At two, $Q$ is odd and $\mathcal E/2\mathcal E$ is the field with four
 elements. Thus $h\equiv1\pmod2$. Put
\[
 s_-=v(h-1)=v_2(P(w)),\qquad
 s_+=v(h+1)=v_2(w^3+\bar w^3).
\]
Both are positive integers. For odd $g$, factoring $h^g-1$ or
$h^g+1$ shows that its valuation is respectively $s_-$ or $s_+$.
For even $g$, the valuations of $h^g-1$ and $h^g+1$ are respectively
$s_-+s_++v_2(g)-1$ and one. Indeed $v_2(h^2-1)\ge2$; every subsequent
doubling increases this valuation by one, and odd multipliers preserve it.
The plus expression is two modulo four. Now
\[
 s_-\log2<\tfrac32\log Q,\qquad
 s_+\log2\le\log2+\tfrac32\log Q.
\]
For even $g$ the $-1$ in the valuation cancels the extra $\log2$.
For odd $g$, $\log2+\tfrac32\log Q\le3\log Q$ since $Q\ge7$.
This proves \eqref{eq:xc-local-two} for both twists.

At three, a primitive pair with norm divisible by three has
$x\equiv y\not\equiv0\pmod3$, so its boundary is a unit. This proves the
claim when $\epsilon=1$. When $\epsilon=0$, $3\nmid Q$ implies
$3\mid P(w)$: modulo three, either a coordinate or their sum is zero.
Thus $v(h-1)=v_3(P(w))+3/2\ge5/2$.
Binomial expansion now gives $v(h^g-1)=v(h-1)+v_3(g)$.
For multiplication by three the last binomial term has valuation $3v(h-1)$,
strictly exceeding $v(h-1)+1$; multipliers prime to three preserve the
leading valuation. Subtracting $v(3\sqrt{-3})=3/2$ proves the exact formula.
The bound on $|P(w)|$ proves \eqref{eq:xc-local-three}.
\end{proof}

\begin{corollary}[A support-free small-prime budget]\label{cor:xc-small}
For $Y\ge3$,
\begin{equation}\label{eq:xc-small}
 \log S_Y(T)\le
 \frac{Y^2+3Y+18}{4}\log Q+\log g.
\end{equation}
In particular, put $Y_g=\max(3,\sqrt g/\log(3g))$. Then
$\log S_{Y_g}(T)=o(\log c)$ as $g\to\infty$, uniformly in the base $w$,
its support, and the twist.
\end{corollary}
\begin{proof}
Sum Theorem~\ref{thm:xc-local}. The contributions of $v_p(g)\log p$ total
at most $\log g$. The two exceptional coefficients sum to $9/2$.
For $n=\lfloor Y\rfloor$, the sum of $p+1$ over supported primes between
five and $Y$ is at most $\sum_{j=1}^n(j+1)=n(n+3)/2$.
This gives \eqref{eq:xc-small}. Dividing by $t$ and using
$t\ge g\log Q/2\ge g\log7/2$, the ratio is at most
\begin{equation}\label{eq:xc-small-ratio}
 \frac{Y_g^2+3Y_g+18}{2g}+\frac{2\log g}{g\log7},
\end{equation}
which tends to zero. This is a uniform elementary bound, not a finite
search assertion or a bound with an unspecified base-dependent constant.
\end{proof}

\section{The resulting restricted abc theorem}
\label{sec:xc-stratum}
Define the explicit error function, with the absolute $C$ from
Theorem~\ref{thm:xc-angle},
\[
 \eta(g)=\frac{2C\log(3g)}g+
 \frac{Y_g^2+3Y_g+18}{2g}+\frac{2\log g}{g\log7}.
\]
It tends to zero.
\begin{theorem}\label{thm:xc-restricted}
Every triple with representation \eqref{eq:xc-content} satisfies
\begin{equation}\label{eq:xc-tail}
 3\log c\le3\log R+\log8+\eta(g)\log c
             +\log E_{3,>Y_g}(T).
\end{equation}
For every $\varepsilon>0$ there is an integer $g_\varepsilon$ such that
\begin{equation}\label{eq:xc-class}
 g\ge g_\varepsilon,\qquad
 E_{3,>Y_g}(T)\le c^{3\varepsilon/(2(1+\varepsilon))}
\end{equation}
imply
\[
 c\le8^{(1+\varepsilon)/3}R^{1+\varepsilon}.
\]
No restriction on $r$ or on an individual $q_i$ is imposed.
\end{theorem}
\begin{proof}
The cubic identity and $N(z)\ge3c^2/4$ give
$T\ge c^3\exp(-C\log Q\log(3g))/8$.
Also $T\le R^3E_3(T)$, and the low-prime contribution to $\log E_3(T)$
is at most $\log S_{Y_g}(T)$. Apply Corollary~\ref{cor:xc-small} and
$\log Q\le2t/g$ to obtain \eqref{eq:xc-tail}.
Choose $g_\varepsilon$ so that $\eta(g)\le3\varepsilon/(2(1+\varepsilon))$
for every $g\ge g_\varepsilon$. Substitution of \eqref{eq:xc-class} gives
$3t/(1+\varepsilon)\le3\log R+\log8$, proving the assertion.
Neither the tail condition nor $g\ge g_\varepsilon$ is asserted on all
triples. In particular, a prime norm has $g=1$ and is not covered by this
argument. No infinitude of the class \eqref{eq:xc-class} is asserted.
\end{proof}

\section{A second compression that allows exponent gcd one}
\label{sec:xc-cluster}
The common exponent is useful but not necessary for compression. For any
integer $m\ge1$, write $e_i=mf_i+r_i$ with $0\le r_i<m$ and put
\[
 W=\prod_i\pi_i^{f_i},\qquad V=\prod_i\pi_i^{r_i},\qquad
 z=u(1+\zeta)^\epsilon VW^m.
\]
Set $Q_W=N(W)$, $Q_V=N(V)$ and
\[
 \mathcal B_m=C_1\max(1,\log Q_V)\max(1,\log Q_W)\log(3m),
\]
where $C_1$ is an effective absolute constant.
\begin{proposition}[Two-block estimates]\label{prop:xc-two-block}
The constant $C_1$ can be chosen such that, uniformly in all the data,
\begin{align}
 |\lambda_z-1|&\ge e^{-\mathcal B_m},\label{eq:xc-block-angle}\\
 v_p(T)&\le p^2\mathcal B_m\quad(p\mid T),\label{eq:xc-block-local}\\
 \log S_Y(T)&\le\mathcal B_mY^3\log Y\quad(Y\ge2).
 \label{eq:xc-block-small}
\end{align}
\end{proposition}
\begin{proof}
The exact multiplicative expression is
$\lambda_z=(-1)^\epsilon(V/\bar V)^3((W/\bar W)^3)^m$.
Apply the complex theorem to its at most three logarithms; a principal-branch
adjustment has coefficient at most $m+2\le3m$.
The two nonconstant heights are bounded by
$3\max(1,\log Q_V)$ and $3\max(1,\log Q_W)$.
For the second assertion use the first inequality of Bugeaud's
Theorem~1.3 \cite{moving:bugeaud}, in degree two, with the same at most three
factors and coefficient bound. One may keep $-1$ and any unit factors with
zero coefficients to have at least two factors. The product is not one.
Since $p\nmid N(z)$, all factors are units at places above $p$, and
$v(\lambda_z-1)=v_p(T)+v(3\sqrt{-3})\ge v_p(T)$.
Enlarging $C_1$ handles both external constants. Summing and using
$\sum_{p\le Y}p^2\log p\le Y^3\log Y$ proves the last assertion.
\end{proof}

\begin{theorem}[Clustered exponents, unrestricted individual norm primes]
\label{thm:xc-cluster}
Fix an integer $J\ge0$, $0<\delta<1$, $B\ge0$ and $\varepsilon>0$.
There is a constant $K_{J,\delta,B,\varepsilon}$ for which
$c\le K_{J,\delta,B,\varepsilon}R^{1+\varepsilon}$ on the following class,
with $t=\log c$ and $c$ sufficiently large:
\begin{equation}\label{eq:xc-cluster-class}
 m\le e_i\le m+J\text{ for all }i,\quad
 m\ge t^{(1+\delta)/2},\quad
 E_{3,>t^{\delta/12}}(T)\le t^B.
\end{equation}
The common gcd of the $e_i$ may be one. There is no separate condition
$q_i\le t^A$ or fixed bound on their number.
\end{theorem}
\begin{proof}
Eventually $m>J$. Then every $f_i=1$. With $S=\sum_i\log q_i$,
$\log Q_W=S\le2t/m$, $\log Q_V\le JS$, and $S\ge\log7$.
Proposition~\ref{prop:xc-two-block} gives
\[
 \mathcal B_m\le C_J(t/m)^2\log(3m)
 \le C'_Jt^{1-\delta}\log t\le t^{1-\delta/2}
\]
past a threshold depending only on $J,\delta$; here $m\le2t/\log7$.
The small-prime bound with $Y=t^{\delta/12}$ is consequently
$O(t^{1-\delta/4}\log t)=o(t)$ uniformly.
The large-prime excess contributes at most $B\log t=o(t)$.
Combine these with $c^3\le8e^{\mathcal B_m}R^3E_3(T)$, and absorb
$\log8+o(t)$ into $3\varepsilon t/(1+\varepsilon)$.
The remaining finite height range is absorbed by one constant.
For example an exponent pattern $(m,m+1)$ has gcd one, so the argument
is not merely the exact common-power case in different notation.
No existence or exhaustion assertion for the class is needed or made.
\end{proof}

\section{Prime auxiliary norm and prescribed high boundary valuations coexist}
\label{sec:xc-ray}
We now attack, rather than assume, a uniform bound on the independent excess
$E_3$. The needed external tool is a quantitative ray-class prime theorem,
not any conjectural height inequality. The ray-class definitions and the
ordinary existence theorem are used in their standard forms
\cite[Chapter~V, Sections~1 and~3; Chapter~VI, Section~4]{MilneCFT}.

\begin{lemma}[A fixed-field least ray-class prime consequence]
\label{lem:xc-ray-prime}
Let $\mathfrak m=3\prod_{p\in S}p^{k_p+1}\mathcal E$, where $S$ is a
nonempty finite set of rational primes at least five and $k_p\ge1$.
Every full ray class modulo $\mathfrak m$ contains a degree-one prime ideal
$\mathfrak q$ satisfying
\[
 N\mathfrak q\le C_2(N\mathfrak m)^{521},
\]
where $C_2$ is effective and absolute. Prime ideals in the class are coprime
to $\mathfrak m$.
\end{lemma}
\begin{proof}
Apply Theorem~3.1 of Thorner--Zaman \cite{XC-ThornerZaman} in the fixed field
$K=\Q(\sqrt{-3})$ of degree two and discriminant three.
Their maximum character-conductor norm, denoted here by $Q_H$, is at most
$N\mathfrak m$. The norm of the product of modulus primes omitted from
the congruence subgroup's conductor is also at most $N\mathfrak m$.
Their threshold is
\[
 3^{694}Q_H^{521}+3^{232}Q_H^{367}2^{580}
 +(12Q_H)^{1/1000}N\mathfrak m_0,
\]
which is at most an absolute multiple of $(N\mathfrak m)^{521}$.
Above the stated conductor threshold the theorem gives a positive lower
bound for the number of degree-one primes in the specified class, hence
at least one.

For completeness, the sufficiently-large-conductor proviso loses only
finitely many moduli in this family. The class number is one, so the full
ray group is $(\mathcal E/\mathfrak m)^\times$ modulo the image of the six
units. For any $p\in S$, the last principal-unit layer
$1+p^{k_p}\mathcal E$ modulo $p^{k_p+1}\mathcal E$ has order $p^2$ and
intersects the unit image trivially. Thus some character is nontrivial on
that layer, and its conductor has a prime-ideal factor of exponent
$k_p+1$. Consequently $Q_H\ge p^{k_p+1}$.
If $Q_H$ is bounded, only finitely many pairs $(p,k_p)$ and hence finitely
many such moduli and classes occur. Existence of a degree-one prime in
each remaining ray class follows from the prime-ideal theorem for ray
classes; alternatively use the corresponding abelian class field and
Chebotarev. The former statement concerns precisely the coprime modulus
classes. Enumerating these finite cases enlarges $C_2$ effectively.
This verifies the hypotheses of the source theorem rather than replacing
its conductor by the much larger discriminant of a varying ray-class field.
\end{proof}

\begin{theorem}[Prime norm with any finite prescribed valuation pattern]
\label{thm:xc-prime-norm}
There is an effective absolute constant $C_{\mathrm{ray}}>0$ such that for every nonempty
finite $S\subset\{p\text{ prime}:p\ge5\}$ and every choice $k_p\ge1$,
there is a primitive positive triple with
\begin{equation}\label{eq:xc-ray-result}
 a^2+ab+b^2=q\text{ prime},\qquad
 v_p(abc)=k_p\quad(p\in S),\qquad
 \prod_{p\in S}p^{k_p}\le c
 \le C_{\mathrm{ray}}\prod_{p\in S}p^{521(k_p+1)}.
\end{equation}
The lower bound holds because all prescribed powers can be placed on one
arm, up to the common unit rotation.
\end{theorem}
\begin{proof}
Put $L=\prod_{p\in S}p^{k_p+1}$ and $A=\prod_{p\in S}p^{k_p}$.
Choose $r\in\mathcal E$ by the ordinary Chinese remainder theorem so that
$r\equiv1\pmod{3\mathcal E}$ and $r\equiv1+A\zeta\pmod{L\mathcal E}$.
It is a unit at the modulus, since its norm is one modulo every $p\in S$.
Apply Lemma~\ref{lem:xc-ray-prime} to the ray class of $(r)$.
Since $\mathcal E$ is a principal ideal domain, equality of ray classes gives
$\mathfrak q=(r\gamma)$ with $\gamma\equiv1\pmod{\mathfrak m}$ locally.
The generator $\pi=r\gamma$ is integral because $(\pi)=\mathfrak q$ is an
integral ideal. Locally at every modulus prime it is congruent to $r$;
since $\pi-r$ is integral, these congruences give $\pi\equiv r\pmod{\mathfrak m}$.
Writing $\pi=x+y\zeta$, we get $x\equiv1$, $y\equiv A\pmod L$.
Thus $v_p(y)=k_p$ and $p\nmid x(x+y)$ for every $p\in S$.

The norm $N(\pi)=q$ is a rational prime. A common divisor of $x,y$ would
have its square divide $q$, so the pair is primitive. Zero boundary would
force a square norm, impossible for a prime. Rotate by a unit to positive
coordinates $a,b$. The absolute values of the three arms are permuted by
this rotation, so all the exact boundary valuations survive, as does the
assertion that $A$ divides a single arm. Hence $c\ge A$.
Finally $N\mathfrak m=9L^2$, and
$c\le2\sqrt q/\sqrt3\le C_{\mathrm{ray}}L^{521}$.
All constants depend only on the fixed quadratic field and the cited
absolute estimate, not on $S$ or the exponents.
\end{proof}

\begin{corollary}[The two uncontrolled regions genuinely overlap]
\label{cor:xc-overlap}
Fix an integer $k\ge4$. There is a family of primitive positive triples
indexed by all rational primes $p\ge5$ with prime auxiliary norm,
$v_p(abc)=k$, and
\[
 p^k\le c\le C_{\mathrm{ray}}p^{521(k+1)}.
\]
For every fixed real $A>0$, eventually $p>(\log c)^A$, and
\begin{equation}\label{eq:xc-large-tail-obstruction}
 \liminf_{p\to\infty}
 \frac{\log E_{3,>(\log c)^A}(abc)}{\log c}
 \ge\frac{k-3}{521(k+1)}>0.
\end{equation}
Consequently the independent estimate
$\log E_{3,>(\log c)^A}(abc)=o(\log c)$ is false even when the auxiliary
norm is prime. The common exponent content of this family is one.
\end{corollary}
\begin{proof}
Use Theorem~\ref{thm:xc-prime-norm} with $S=\{p\}$.
The upper height bound gives $\log c\le\log C_{\mathrm{ray}}+521(k+1)\log p$.
Every fixed power of this logarithm is eventually less than $p$.
The specified prime then contributes $p^{k-3}$ to the large-prime excess,
giving \eqref{eq:xc-large-tail-obstruction} on division.
Since $N(z)$ is prime, \eqref{eq:xc-factorization} has a single exponent one.
The family is unbounded because $c\ge p^k$.
\end{proof}

This is not an abc counterfamily. No upper bound for its full radical $R$
strong enough to violate abc has been proved. The complementary credit
$C_3(T)$ and the balance of the three arms remain relevant. In particular,
this corollary does not contradict the previous restricted theorem, whose
small-norm support assumptions fail on prime norms of size comparable to
$c^2$. It rules out extending that theorem merely by asserting its separate
large-prime excess hypothesis everywhere.

\section{An intrinsic angular obstruction and a modular cross-check}
\label{sec:xc-audit}
\begin{proposition}[A sharp elementary universal angular scale]
\label{prop:xc-angular-cap}
For every primitive positive triple,
\[
 |\lambda_z-1|\ge\frac{3\sqrt3}{2c},\qquad
 \max(0,-\log|\lambda_z-1|)\le\log c.
\]
Thus the preceding $O(\log c\log\log c)$ upper penalty at prime norm can
already be replaced by an $O(\log c)$ penalty without any logarithmic-form
input. In the exact-power class one may take the minimum of $\log c$ and
$C\log Q\log(3g)$ as a valid nonnegative angular penalty.
\end{proposition}
\begin{proof}
Since $(a-1)(b-1)\ge0$, $ab\ge c-1\ge c/2$. Hence $T\ge c^2/2$.
Use $|\lambda_z-1|=3\sqrt3T/M^{3/2}$ and $M\le c^2$. Since
$3\sqrt3/2>1$, the stated logarithmic consequence follows, including when
the absolute value is greater than one. Taking the smaller of two valid
upper penalties preserves the lower bound for the absolute value.
\end{proof}

\begin{proposition}[Angular smallness is not a universal target]
\label{prop:xc-imbalance}
On the genuine primitive family $(1,c-1,c)$,
\[
 \left|\left(\frac z{\bar z}\right)^3-1\right|
 =\frac{3\sqrt3\,c(c-1)}{(c^2-c+1)^{3/2}},\qquad
 -\log|\lambda_z-1|=\log c-\log(3\sqrt3)+o(1).
\]
In particular no universally sublinear angular penalty can be valid on all
triples, regardless of the chosen decomposition of $z$.
\end{proposition}
\begin{proof}
Substitute $T=c(c-1)$ and $N(z)=c^2-c+1$ in the cubic identity and divide
numerator and denominator by $c^3$. The residual factor tends to one.
Any $A(z)$ with $|\lambda_z-1|\ge e^{-A(z)}$ must therefore satisfy
$A(z)\ge\log c+O(1)$ on this family.
\end{proof}

Thus two separate universal replacements for the previous sufficient class
are now excluded: unconditionally making the angular penalty sublinear,
and unconditionally making the large-prime excess sublinear. Their
correlated, credit-preserving alternatives remain open, not refuted.

We also tested the modular direction against these proposed replacements.
Pasten's established estimate \cite[Theorem~2.5]{XC-Pasten} states that,
for each $\rho>0$,
\[
 \prod_{p\mid T}v_p(T)\le\kappa_\rho R^{8/3+\rho}.
\]
It implies, for $V>1$,
\begin{equation}\label{eq:xc-pasten-count}
 \#\{p\mid T:v_p(T)\ge V\}
 \le\frac{\log\kappa_\rho+(8/3+\rho)\log R}{\log V}.
\end{equation}
The proof is simply to take logarithms and retain these factors.
The source controls $\sum\log v_p(T)$, not
$\sum(v_p(T)-3)_+\log p$. We do not replace one by the other.
Corollary~\ref{cor:xc-overlap}, with a fixed $k\ge4$ and $p\to\infty$,
shows concretely why excluding only extremely large exponents cannot by
itself exclude a large weighted excess: this obstruction uses a fixed
exponent. No contradiction with Pasten's theorem is asserted.

\section{Finite certificates, formalization scope, and the surviving graph}
\label{sec:xc-validation}
The supplied standard-library Python verifier checks the local inequalities
without logarithmic rounding. For $p\ge5$ it verifies
\[
 p^{2\max(v_p(T)-v_p(g),0)}\le Q^{p+1};
\]
the exceptional-prime and aggregate estimates are likewise checked by
integer powers. The aggregate comparison used in the replay is
$S_Y(T)^4\le g^4Q^{Y^2+3Y+18}$ at integer $Y$.
The test uses positive primitive bases with coordinates at most 24,
$3\nmid Q$, $1\le g\le12$, both twists, and all primes at most 101.
It checks 6,528 triples and 169,728 local prime conditions, plus
52,224 aggregate inequalities. An independently implemented integer matrix
power verifies 3,264 pair powers. Forty-two adjacent-exponent examples have
exponent gcd one. None of these counts proves a universal theorem.

Five concrete prime-norm examples have the shape
$(1,p^4-1,p^4)$ for $p=5,13,31,107,263$.
For example the last has
\[
 (a,b,c)=(1,4784350560,4784350561),\qquad
 N(z)=22890010285756664161\text{ prime}.
\]
Their prime claims are checked by recursive full-$n-1$ Lucas certificates,
not by probable-prime tests. The certificate graph has 21 prime nodes and
81 modular conditions. To justify this test, if $a^{n-1}=1\pmod n$ and
$\gcd(a^{(n-1)/\ell}-1,n)=1$ for every certified prime $\ell\mid n-1$,
then for every prime divisor $r$ of $n$ the order of $a$ modulo $r$ contains
the full $\ell$-part of $n-1$. Thus $n-1\mid r-1$, forcing $r=n$.
This proves the criterion recursively from the prime two. The last example
has exact radical $R=78642762330>c$, despite $v_{263}(T)=4$ and prime
auxiliary norm; thus a large selected excess can coexist with ample radical.
These five
examples are not substituted for the ray-class infinite-family proof.

A separate Lean companion records elementary polynomial and divisibility
statements, and integer forms of selected bounds. It is expressly marked \emph{uncompiled}:
there is no Lean executable in this execution and no accessible remote
write/CI-trigger action was provided by the current connector.
No number-theoretic existence result, local order theorem, logarithmic-form
estimate, ray-class theorem or standard abc statement is thereby claimed
kernel-checked. The earlier 70-declaration validation records are retained
as historical records and are not represented as validation of new code.

\begin{center}\small
\begin{tabular}{p{.21\linewidth}p{.34\linewidth}p{.37\linewidth}}
\toprule
Line & Result established here & Exact remaining issue\\\midrule
Exact exponent compression & Support-free angular and elementary local bounds
& Large-prime tail and small exponent content\\
Near-common exponents & Three-log compression even at gcd one
& Coverage beyond clustered exponents and large-prime tail\\
Ray-class construction & Prime norm and deep large boundary primes coexist
& Full radical/credit correlation, not just excess\\
Modular multiplicity & Literal consequence \eqref{eq:xc-pasten-count}
& Conversion from log exponent count to weighted depth\\
IUT and other geometry & No new source comparison asserted
& Previous all-place and uniform height obligations\\
\bottomrule
\end{tabular}
\end{center}

The positive dependency chain is
\[
 \text{actual power extraction}+
 \text{two-log estimate}+
 \text{explicit local lifting}
 \ \Longrightarrow\ \eqref{eq:xc-tail}.
\]
Its final implication to abc is proved only under \eqref{eq:xc-class}.
The second positive chain uses the three-log and $p$-adic input only on
\eqref{eq:xc-cluster-class}. The adversarial chain is
\[
 \text{least ray-class prime theorem}+\text{CRT}
 \ \Longrightarrow\ \eqref{eq:xc-ray-result}
 \ \Longrightarrow\ \eqref{eq:xc-large-tail-obstruction}.
\]
Every external input is an established theorem, but no chain removes all
the remaining restrictions. In particular no unconditional term of type
\texttt{ABCConjecture} is constructed. Claims about the local checks,
source integration and remote publication are recorded separately from
mathematical correctness. No independent autonomous subagents or external
peer review were available in this execution.
% END EXPONENT-COMPRESSION SUPPLEMENT 2026-09-06
